% ==========================================================
% 5. FORMULACIÓN DE HIPÓTESIS (Fase 3, Tema 3.6)
% ==========================================================

\section{Formulación de hipótesis}

Por su alcance correlacional, el estudio formula hipótesis correlacionales con la
dirección esperada, una por cada relación que los antecedentes y la matriz de
correlación (Tabla~2) sustentan, cada una con su hipótesis nula.

\medskip
\noindent\textbf{Hipótesis 1: exposición a pantallas y fatiga visual.}\\
H$_{1a}$: a mayor tiempo diario de exposición a pantallas durante las jornadas de
estudio, mayor frecuencia de fatiga visual en los estudiantes de Ingeniería de Sistemas
y Computación de la UPTC [5].\\
H$_{0a}$: no existe relación entre el tiempo diario de exposición a pantallas y la
frecuencia de fatiga visual.

\medskip
\noindent\textbf{Hipótesis 2: exposición a pantallas y molestias osteomusculares.}\\
H$_{1b}$: a mayor tiempo diario de exposición a pantallas, mayor frecuencia de
molestias osteomusculares [7].\\
H$_{0b}$: no existe relación entre el tiempo diario de exposición a pantallas y la
frecuencia de molestias osteomusculares.

\medskip
\noindent\textbf{Hipótesis 3: hábitos posturales y molestias osteomusculares.}\\
H$_{1c}$: los estudiantes con hábitos posturales más inadecuados durante el estudio
presentan mayor frecuencia de molestias osteomusculares [1], [6].\\
H$_{0c}$: no existe relación entre los hábitos posturales y la frecuencia de molestias
osteomusculares.

\medskip
\noindent No se formula una hipótesis entre hábitos posturales y fatiga visual porque la
matriz de correlación muestra una relación débil o inexistente entre ambas (Tabla~2,
causa C4 con S1 y S2). Cada hipótesis es falsable: una ausencia de relación, o una
relación en sentido contrario, la refutaría.
